% Generated draft: requires mathematical and source review.
\documentclass[11pt]{article}
\usepackage[margin=25mm]{geometry}
\usepackage{fontspec}
\setmainfont{Latin Modern Roman}
\usepackage{amsmath,amssymb,mathtools,unicode-math}
\setmathfont{Latin Modern Math}
\usepackage{enumitem,needspace,xurl,hyperref}
\setlength{\parindent}{0pt}
\setlength{\parskip}{5pt}
\setlength{\emergencystretch}{3em}
\newcommand{\N}{\mathbb{N}}
\newcommand{\Z}{\mathbb{Z}}
\newcommand{\Q}{\mathbb{Q}}
\newcommand{\R}{\mathbb{R}}
\newcommand{\C}{\mathbb{C}}
\newcommand{\F}{\mathbb{F}}
\newcommand{\A}{\mathbb{A}}
\newcommand{\PP}{\mathbb{P}}
\newcommand{\E}{\mathbb{E}}
\newcommand{\eps}{\varepsilon}
\newcommand{\dd}{\,\mathrm{d}}
\begin{document}
\section*{UnsolvedMath ID: 800719}
\textbf{Source problem code:} catalogue entry
\subsection{Definitions and mathematical statement}
Let \(F\) be a field. Interpret the division algebras in the source as finite-dimensional central division algebras \(D_1,D_2\) over the same field \(F\), so that their Brauer classes \([D_1],[D_2]\in\operatorname{Br}(F)\) are defined. This common-base, centrality, and finite-dimensionality interpretation is necessary for the usual Brauer-group formulation but is not explicit in the extracted statement.

For a field extension \(L/F\), say that \(L\) splits \(D_i\) if \(D_i\otimes_F L\) is isomorphic to a full matrix algebra over \(L\), equivalently if the restriction of \([D_i]\) to \(\operatorname{Br}(L)\) is zero. For a specified class \(\mathcal C(F)\) of extensions of \(F\), define
\[
\operatorname{Spl}_{\mathcal C}(D_i)=\{L/F\in\mathcal C(F):L\text{ splits }D_i\}.
\]
The subgroup generated by a Brauer class is
\[
\langle[D_i]\rangle=\{m[D_i]:m\in\Z\}\subseteq\operatorname{Br}(F).
\]

The question is whether, for every field \(F\) and every such pair \(D_1,D_2\), the implication
\[
\operatorname{Spl}_{\mathcal C}(D_1)=\operatorname{Spl}_{\mathcal C}(D_2)
\quad\Longrightarrow\quad
\langle[D_1]\rangle=\langle[D_2]\rangle
\]
holds for the intended class \(\mathcal C(F)\). The source specifically suggests the following two choices, to be considered separately:
\[
\mathcal C_{\mathrm{fin}}(F)=\{L/F:[L:F]<\infty\},
\qquad
\mathcal C_{\leq 1}(F)=\{L/F:\operatorname{trdeg}_F L\leq 1\}.
\]
The source does not fix a single class or specify restrictions on the base field. The displayed classes follow the literal wording: no separability requirement is imposed on finite extensions, and no finite-generation requirement is imposed on extensions of transcendence degree at most one. Any intended additional restrictions remain unresolved.
\subsection{Short English statement}
For two central division algebras over a common field, ask whether having exactly the same splitting extensions within a specified class forces their Brauer classes to generate the same cyclic subgroup. The source suggests finite extensions and extensions of transcendence degree at most one, without fixing additional hypotheses.
\subsection{Sources}
\begin{itemize}
\item Local UnsolvedMath registry, permanent ID \(800719\), source problem code \texttt{catalogue entry}, titled \emph{Open Problems on Central Simple Algebras --- Problem 33}. The supplied registry labels the problem open; this label is not an independently verified mathematical or literature assessment.
\item Auel, Brussel, Garibaldi, and Vishne, \emph{Open Problems on Central Simple Algebras} (2010), Problem 33. Supplied source URL: \url{https://arxiv.org/abs/1006.3304}. The registry attributes the problem to these authors.
\item Supplied provenance records access on 2026-10-16 and extraction by \texttt{arxiv\_latex}. These are supplied provenance details; the source has not been independently accessed or verified for this draft.
\end{itemize}
\end{document}
